\documentclass[aps,preprint]{revtex4}%
\usepackage{amssymb}
\usepackage{amsfonts}%
\usepackage{amsmath}%
\setcounter{MaxMatrixCols}{30}%
\usepackage{graphicx}
%TCIDATA{OutputFilter=latex2.dll}
%TCIDATA{Version=5.50.0.2960}
%TCIDATA{CSTFile=revtex4.cst}
%TCIDATA{Created=Friday, July 12, 2019 20:12:44}
%TCIDATA{LastRevised=Thursday, July 18, 2019 12:34:54}
%TCIDATA{<META NAME="GraphicsSave" CONTENT="32">}
%TCIDATA{<META NAME="SaveForMode" CONTENT="1">}
%TCIDATA{BibliographyScheme=Manual}
%TCIDATA{<META NAME="DocumentShell" CONTENT="Articles\SW\REVTeX 4">}
%BeginMSIPreambleData
\providecommand{\U}[1]{\protect\rule{.1in}{.1in}}
%EndMSIPreambleData
\newtheorem{theorem}{Theorem}
\newtheorem{acknowledgement}[theorem]{Acknowledgement}
\newtheorem{algorithm}[theorem]{Algorithm}
\newtheorem{axiom}[theorem]{Axiom}
\newtheorem{claim}[theorem]{Claim}
\newtheorem{conclusion}[theorem]{Conclusion}
\newtheorem{condition}[theorem]{Condition}
\newtheorem{conjecture}[theorem]{Conjecture}
\newtheorem{corollary}[theorem]{Corollary}
\newtheorem{criterion}[theorem]{Criterion}
\newtheorem{definition}[theorem]{Definition}
\newtheorem{example}[theorem]{Example}
\newtheorem{exercise}[theorem]{Exercise}
\newtheorem{lemma}[theorem]{Lemma}
\newtheorem{notation}[theorem]{Notation}
\newtheorem{problem}[theorem]{Problem}
\newtheorem{proposition}[theorem]{Proposition}
\newtheorem{remark}[theorem]{Remark}
\newtheorem{solution}[theorem]{Solution}
\newtheorem{summary}[theorem]{Summary}
\newenvironment{proof}[1][Proof]{\noindent\textbf{#1.} }{\ \rule{0.5em}{0.5em}}
\begin{document}
\preprint{HEP/123-qed}
\title[Short title for running header]{REV\TeX4}
\author{First Author}
\affiliation{My Institution}
\author{Second Author}
\affiliation{My Institution}
\author{Third Author}
\affiliation{Other Institution}
\keywords{one two three}
\pacs{PACS number}

\begin{abstract}
Shell document for REV\TeX{} 4.

\end{abstract}
\volumeyear{year}
\volumenumber{number}
\issuenumber{number}
\eid{identifier}
\date[Date text]{date}
\received[Received text]{date}

\revised[Revised text]{date}

\accepted[Accepted text]{date}

\published[Published text]{date}

\startpage{101}
\endpage{102}
\maketitle
\tableofcontents


\section{Orthogonal Transformations }

Consider a vector space $V$ over the field $\mathbb{F}$ equipped with a
positive definite inner product%
\begin{equation}
\left\langle \left.  {}\right\vert \right\rangle :V\times V\rightarrow
\mathbb{F}%
\end{equation}
and bases of vectors $\left\{  \left.  \varphi_{j}\right\vert 1\leq j\leq
r\right\}  $ and $\left\{  \left.  \psi_{j}\right\vert 1\leq j\leq r\right\}
$ that are orthogonal wrt $\left\langle \left.  {}\right\vert \right\rangle $
i.e.
\begin{equation}
\left\langle \left.  \psi_{j}\right\vert \psi_{k}\right\rangle =\delta
_{jk};1\leq j,k\leq r
\end{equation}
and%
\begin{equation}
\left\langle \left.  \varphi_{j}\right\vert \varphi_{k}\right\rangle
=\delta_{jk};1\leq j,k\leq r.
\end{equation}


If the bases are normalized then the transition operators%
\begin{equation}
\left\vert \varphi_{j}\right\rangle \left\langle \psi_{j}\right\vert :\psi
_{k}\rightarrow\varphi_{k};1\leq k\leq r
\end{equation}
are partial isometries and their sum
\begin{equation}%
%TCIMACRO{\tsum \limits_{1\leq j\leq r}}%
%BeginExpansion
{\textstyle\sum\limits_{1\leq j\leq r}}
%EndExpansion
\left\vert \varphi_{j}\right\rangle \left\langle \psi_{j}\right\vert :\psi
_{k}\rightarrow\varphi_{k};1\leq k\leq r
\end{equation}
is an isometry. All partial isometries and isometries are given by
\begin{equation}
\left\vert O_{1}\varphi_{j}\right\rangle \left\langle O_{2}\psi_{j}\right\vert
;1\leq j\leq r,O_{1},O_{2}\in O\left(  r,\mathbb{F}\right)
\end{equation}
and
\begin{equation}%
%TCIMACRO{\tsum \limits_{1\leq j\leq r}}%
%BeginExpansion
{\textstyle\sum\limits_{1\leq j\leq r}}
%EndExpansion
\left\vert O_{1}\varphi_{j}\right\rangle \left\langle O_{2}\psi_{j}\right\vert
,O_{1},O_{2}\in O\left(  r,\mathbb{F}\right)  .
\end{equation}


Any diagonal operator can be expressed as
\begin{equation}%
%TCIMACRO{\tsum \limits_{1\leq j\leq r}}%
%BeginExpansion
{\textstyle\sum\limits_{1\leq j\leq r}}
%EndExpansion
\lambda_{j}\left\vert \varphi_{j}\right\rangle \left\langle \varphi
_{j}\right\vert
\end{equation}
so that
\begin{equation}%
%TCIMACRO{\tsum \limits_{1\leq j\leq r}}%
%BeginExpansion
{\textstyle\sum\limits_{1\leq j\leq r}}
%EndExpansion
\lambda_{j}\left\vert \varphi_{j}\right\rangle \left\langle \varphi
_{j}\right\vert :\varphi_{k}\rightarrow\lambda_{k}\varphi_{k}%
\end{equation}
if
\begin{equation}
\left\langle \left.  \varphi_{j}\right\vert \varphi_{k}\right\rangle
=\delta_{jk}%
\end{equation}


A transformation $T\in GL\left(  V,\mathbb{F}\right)  $ has the effect
\begin{equation}
T:v_{k}\rightarrow w_{k},
\end{equation}
where if $\left\{  \left.  v_{j}\right\vert 1\leq j\leq r\right\}  $ are
linearly independent then $\left\{  \left.  w_{j}\right\vert 1\leq j\leq
r\right\}  =\left\{  T\left(  w_{j}\right)  \right\}  $ are linearly
independent.
\begin{equation}
Oe_{j}\bot Oe_{k}%
\end{equation}%
\begin{align}
&
%TCIMACRO{\tsum \limits_{1\leq j\leq r}}%
%BeginExpansion
{\textstyle\sum\limits_{1\leq j\leq r}}
%EndExpansion
\left\vert Oe_{j}\right\rangle \left\langle e_{j}\right\vert \\%
%TCIMACRO{\tsum \limits_{1\leq j\leq r}}%
%BeginExpansion
{\textstyle\sum\limits_{1\leq j\leq r}}
%EndExpansion
\left\vert Of_{j}\right\rangle \left\langle f_{j}\right\vert  & =%
%TCIMACRO{\tsum \limits_{1\leq j\leq r}}%
%BeginExpansion
{\textstyle\sum\limits_{1\leq j\leq r}}
%EndExpansion
\left\vert OO_{1}e_{j}\right\rangle \left\langle O_{1}e_{j}\right\vert
\end{align}%
\begin{equation}
A=\sum_{1\leq j\leq s}i\alpha_{j}\left\{  \left\vert v_{j}\right\rangle
\left\langle v_{j}\right\vert -\left\vert v_{j}^{\ast}\right\rangle
\left\langle v_{j}^{\ast}\right\vert \right\}  ;\qquad\alpha_{j}\in
\mathbb{R}^{+}\label{sa}%
\end{equation}


\begin{theorem}
Any orthogonal transformation $O$ can be expanded
\begin{equation}
O=\sum\left\vert f_{j}\right\rangle \left\langle e_{j}\right\vert
\end{equation}

\end{theorem}

\begin{proof}%
\begin{equation}
O=\exp\left\{  A\right\}  =\exp\left\{  \sum_{1\leq j\leq s}i\alpha
_{j}\left\{  \left\vert v_{j}\right\rangle \left\langle v_{j}\right\vert
-\left\vert v_{j}^{\ast}\right\rangle \left\langle v_{j}^{\ast}\right\vert
\right\}  \right\}
\end{equation}

\end{proof}


\end{document}